%%%%%%%%%%%%%%%%%%%%%%%%%%%%%%%%%%%%%%%%%%%%%%%%%%%%%%%
% CHAPTER: Hard-TireNet
%%%%%%%%%%%%%%%%%%%%%%%%%%%%%%%%%%%%%%%%%%%%%%%%%%%%%%%
\chapter{Hard-TireNet: Hard-Constrained Combined-Slip Tire-Force Surrogates}
\label{Tire:chap:Hard-TireNet}
This chapter develops a physical application of HardNet++ for combined-slip
tire-force estimation in autonomous vehicles. The proposed framework,
referred to as Hard-TireNet, learns the nonlinear longitudinal and lateral
tire forces while strictly satisfying a prescribed load- and friction-dependent
admissible set. In prior work, vanilla HardNet was applied
after approximating the nonlinear friction ellipse with conservative affine constraints.
In contrast, this work leverages the general formulation of HardNet++~\cite{Goertzen:2026aa}
to directly enforce the nonlinear friction ellipse without polyhedral over-approximation.
The force-envelope limits are assumed to be known a priori.
The remainder of this chapter formulates the combined-slip estimation
problem in Section~\ref{Tire:sec:problem:formulation}, defines the physical
admissible set in Section~\ref{Tire:sec:physical:constraints}, introduces
the Hard-TireNet architecture in Section~\ref{Tire:sec:architecture},
describes training and online adaptation in Section~\ref{Tire:sec:training},
and presents the validation protocol in Section~\ref{Tire:sec:validation}.
The chapter is concluded in Section~\ref{Tire:sec:conclusion}.

%%%%%%%%%%%%%%%%%%%%%%%%%%%%%%%%%%%%%%%%%%%%%%%%%%%%%%%
% PROBLEM FORMULATION
%%%%%%%%%%%%%%%%%%%%%%%%%%%%%%%%%%%%%%%%%%%%%%%%%%%%%%%
\section{Problem Formulation}
\label{Tire:sec:problem:formulation}

Under combined longitudinal and lateral slip, the tire-force vector is
\begin{equation}
  \label{Tire:eq:force:vector}
  \mv{F}
  :=
  (
    F_x,
    F_y
  )^\top
  ,
\end{equation}
where $F_x$ and $F_y$ are the longitudinal and lateral tire forces,
respectively.  Under steady-state conditions, the tire forces are nonlinearly
dependent on the  longitudinal slip ratio $\kappa$, slip angle $\alpha$, 
vertical load $F_z$, road-friction coefficient $\mu$, and tire-specific 
parameters.  We write this relationship as
\begin{equation}
  \label{Tire:eq:force:map}
  \mv{F}
  =
  \mv{f}_{\mathrm{tire}}
  \left(
    \mv{x}; \mv{\eta}
  \right)
  ,
\end{equation}
where $\mv{x}:=(\kappa,\alpha,F_z,\mu)^\top$ and $\mv{\eta}$ collects unknown or
slowly varying tire and road parameters.

In autonomous vehicle control and state estimation, evaluating the exact
tire-force map $\mv{f}_{\mathrm{tire}}$ in real time is challenging due to
modeling complexity and environmental uncertainties. To address this, the
objective of Hard-TireNet is to construct a neural surrogate
$\widehat{\mv{F}} := (\widehat{F}_x,\widehat{F}_y)^\top$ that accurately
approximates the nonlinear force map while guaranteeing that the estimated
forces strictly obey the physical friction envelope. This surrogate problem
is formulated as

\begin{equation}
  \label{Tire:eq:objective}
  \widehat{\mv{F}}
  \approx
  \mv{F},
  \qquad
  \text{subject to}
  \qquad
  \widehat{\mv{F}}
  \in
  \mathcal{C}
  \left(
    \mv{x};\mv{\eta}
  \right),
\end{equation}
where $\mathcal{C}(\mv{x};\mv{\eta}) \subset \R^2$ denotes the state-dependent physical
admissible set governed by the combined-slip friction ellipse. Unlike standard
unconstrained or soft-penalized neural networks that may predict unphysical
forces under distribution shifts or sparse data, the condition in
\eqref{Tire:eq:objective} is enforced as a hard constraint. The mathematical
structure of $\mathcal{C}(\mv{x};\mv{\eta})$ is formally defined in Section~\ref{Tire:sec:physical:constraints}.

\subsection{Assumptions}
\label{Tire:sec:assumptions}

The initial study is based on the following assumptions.

\begin{assum}
  \label{Tire:assum:context}
  The vertical load and road-friction coefficient are available to the
  surrogate, either from measurement or from separate vehicle-state
  observers.  They satisfy $F_z\in\R_{>0}$ and
  $\mu\in[\mu_{\min},\mu_{\max}]$ on the operating region.
\end{assum}

\begin{assum}
  \label{Tire:assum:compact}
  The combined-slip input belongs to a compact set
  $\Omega_x\subset\R^4$, and the reference force map
  $\mv{f}_{\mathrm{tire}}$ is continuous on $\Omega_x$.
\end{assum}

\begin{assum}
  \label{Tire:assum:envelope}
  The combined-slip admissible set $\mathcal{C}(\mv{x};\mv{\eta})$ is nonempty on $\Omega_x$.
  Specifically, the force-envelope boundary functions $F_{x,\lim}(\mv{x};\mv{\eta})$ and
  $F_{y,\lim}(\mv{x};\mv{\eta})$ (known a priori)
  are strictly positive and continuous on the operating domain.
\end{assum}
Assumption~\ref{Tire:assum:envelope} ensures that the physical envelope has a
well-defined interior with non-zero semi-axes. This is critical for HardNet++,
as the differentiable projection layer can only project onto a feasible, non-empty
constraint set; an inconsistent or degenerate constraint boundary cannot be repaired
by the network.

%%%%%%%%%%%%%%%%%%%%%%%%%%%%%%%%%%%%%%%%%%%%%%%%%%%%%%%
% PHYSICAL CONSTRAINTS
%%%%%%%%%%%%%%%%%%%%%%%%%%%%%%%%%%%%%%%%%%%%%%%%%%%%%%%
\section{Combined-Slip Physical Constraint}
\label{Tire:sec:physical:constraints}

The tire forces $\mv{F}=(F_x,F_y)^\top$ are constrained by the friction ellipse
\begin{equation}
  \label{Tire:eq:friction:ellipse}
  \left(
  \frac{F_x}{F_{x,\lim}(\mv{x};\mv{\eta})}
  \right)
  ^2
  +
  \left(
  \frac{F_y}{F_{y,\lim}(\mv{x};\mv{\eta})}
  \right)
  ^2
  \leq
  1
  ,
\end{equation}
where $F_{x,\lim}(\mv{x};\mv{\eta})$ and $F_{y,\lim}(\mv{x};\mv{\eta})$ are the
longitudinal and lateral force limits, respectively.  They depend on the
vertical load, road friction, and tire-road parameters.  Define the scalar
constraint function
\begin{equation}
  \label{Tire:eq:friction:ellipse:rewritten}
  c
  \left(
    \mv{F};\mv{x},\mv{\eta}
  \right)
  :=
  \left(
  \frac{F_x}{F_{x,\lim}(\mv{x};\mv{\eta})}
  \right)^2
  +
  \left(
  \frac{F_y}{F_{y,\lim}(\mv{x};\mv{\eta})}
  \right)^2 
  -1
  .
\end{equation}

The combined-slip force constraint set introduced in
\eqref{Tire:eq:objective} is therefore represented as
\begin{equation}
  \label{Tire:eq:general:set}
  \mathcal{C}
  \left(
    \mv{x};\mv{\eta}
  \right)
  :=
  \left\{
    \mv{F}\in\R^2
    \mid
    c
    \left(
      \mv{F};\mv{x},\mv{\eta}
    \right)
    \leq
    0
  \right\}
  .
\end{equation}

%%%%%%%%%%%%%%%%%%%%%%%%%%%%%%%%%%%%%%%%%%%%%%%%%%%%%%%
% PROPOSED METHOD
%%%%%%%%%%%%%%%%%%%%%%%%%%%%%%%%%%%%%%%%%%%%%%%%%%%%%%%
\section{Proposed Method}
\label{Tire:sec:architecture}
The proposed architecture is shown in Fig.~\ref{Tire:fig:architecture}.  It
contains a force-prediction backbone and a
HardNet++ differentiable projection layer.  The backbone produces an
unconstrained candidate force, while the projection layer maps that candidate
to the current combined-slip admissible set before it is used for training,
  estimation, or control.  Throughout this section, $\mv{F}$ denotes the true
force, $\widehat{\mv{F}}_u$ the raw backbone output, and
$\widehat{\mv{F}}$ the constrained output.

\begin{figure}[t]
  \centering
  \resizebox{0.8\textwidth}{!}{%
    \begin{tikzpicture}[
    x=1cm,
    y=1cm,
    >=Latex,
    font=\small,
    block/.style={
        draw=black!90,
        rounded corners=2pt,
        line width=0.8pt,
        align=center,
        inner sep=5pt
    },
    inputbox/.style={
        block,
        fill=blue!6,
        minimum width=3.9cm,
        minimum height=2.35cm
    },
    sharedbox/.style={
        block,
        fill=blue!6,
        minimum width=6.3cm,
        minimum height=4.8cm
    },
    baseline/.style={
        block,
        fill=gray!10,
        text width=4.6cm,
        minimum height=1.35cm
    },
    softouter/.style={
        block,
        fill=yellow!14,
        text width=4.6cm,
        minimum height=2.45cm
    },
    hardouter/.style={
        block,
        fill=green!16,
        text width=4.6cm,
        minimum height=3.15cm
    },
    innerlayer/.style={
        draw=black!80,
        rounded corners=2pt,
        line width=0.65pt,
        fill=white,
        align=center,
        inner sep=4pt,
        text width=3.55cm
    },
    context/.style={
        block,
        fill=orange!8,
        draw=orange!85!black,
        minimum width=3.9cm,
        minimum height=1.35cm
    },
    chronoPlaceholder/.style={
        draw=black!55,
        dashed,
        rounded corners=2pt,
        line width=0.7pt,
        fill=gray!3,
        align=center,
        minimum width=2.8cm,
        minimum height=1.15cm,
        inner sep=3pt
    },
    conn/.style={
        -{Latex[length=2.4mm,width=1.7mm]},
        line width=0.8pt,
        draw=black!82
    },
    dashedconn/.style={
        -{Latex[length=2.2mm,width=1.5mm]},
        dashed,
        line width=0.75pt,
        draw=black!70
    },
    dense/.style={
        draw=black!55,
        line width=0.50pt
    },
    neuron/.style={
        circle,
        draw=black!80,
        fill=blue!10,
        line width=0.7pt,
        minimum size=6.5mm,
        inner sep=0pt
    }
]

% ================================================================
% Task inputs
% ================================================================
\node[
    inputbox,
    very thick,
    node font=\small\bfseries
] (input) at (-0.32,0.02) {
    Task inputs\\
    $\text{lateral: }(\alpha,F_z,\mu)$\\
    $\text{combined: }(\kappa,\alpha,F_z,\mu)$
};

% ================================================================
% Project Chrono data-generation source
% ================================================================

\def\chronoXShift{0.0cm}   % + right, - left
\def\chronoYShift{1.55cm}  % + up,    - down

\node[
    chronoPlaceholder,
    minimum width=115pt,
    minimum height=48pt,
    xshift=\chronoXShift,
    yshift=\chronoYShift
] (chrono) at (input.north) {
    \includegraphics[width=2.75cm]{figs/Hard-TireNet/pC.png}\\[-1pt]
    {\scriptsize\bfseries Data generation}
};

\draw[
    dashed,
    -{Latex[length=2.0mm,width=1.4mm]},
    line width=0.65pt,
    draw=black!55
]
(chrono.south) -- (input.north);

% ================================================================
% Physical context
% ================================================================
\node[
    context,
    very thick,
    node font=\small\bfseries,
    minimum height=62pt
] (phys) at (-0.25,-3.87) {
    Physical context\\
    $(F_z,\mu)$
};

% ================================================================
% Common MLP architecture
% ================================================================
\node[
    sharedbox,
    very thick,
    minimum height=168pt
] (shared) at (5.85,0) {};

% Movable shared-title position
\def\sharedTitleXShift{0.0cm}   % + right, - left
\def\sharedTitleYShift{-0.7cm}  % + up,    - down

\node[
    node font=\large\bfseries,
    align=center,
    xshift=\sharedTitleXShift,
    yshift=\sharedTitleYShift
] at (shared.north) {
    Common fully connected\\
    MLP architecture
};

% ------------------------------------------------
% Neural-network nodes
% ------------------------------------------------
\begin{scope}[scale=1.17, transform shape]

\node[neuron] (a1) at ($(shared.center)+(-2.00, 1.25)$) {};
\node[neuron] (a2) at ($(shared.center)+(-2.00, 0.42)$) {};
\node[neuron] (a3) at ($(shared.center)+(-2.00,-0.42)$) {};
\node[neuron] (a4) at ($(shared.center)+(-2.00,-1.25)$) {};

\node[neuron] (b1) at ($(shared.center)+(-0.72, 0.95)$) {};
\node[neuron] (b2) at ($(shared.center)+(-0.72, 0.00)$) {};
\node[neuron] (b3) at ($(shared.center)+(-0.72,-0.95)$) {};

\node[neuron] (c1) at ($(shared.center)+(0.75, 0.95)$) {};
\node[neuron] (c2) at ($(shared.center)+(0.75, 0.00)$) {};
\node[neuron] (c3) at ($(shared.center)+(0.75,-0.95)$) {};

\node[neuron] (d1) at ($(shared.center)+(1.95, 0.68)$) {};
\node[neuron] (d2) at ($(shared.center)+(1.95,-0.68)$) {};

\end{scope}

% ------------------------------------------------
% Dense connections
% ------------------------------------------------
\foreach \s in {a1,a2,a3,a4}{
    \foreach \t in {b1,b2,b3}{
        \draw[dense] (\s) -- (\t);
    }
}

\foreach \s in {b1,b2,b3}{
    \foreach \t in {c1,c2,c3}{
        \draw[dense] (\s) -- (\t);
    }
}

\foreach \s in {c1,c2,c3}{
    \foreach \t in {d1,d2}{
        \draw[dense] (\s) -- (\t);
    }
}

% ================================================================
% Baseline branch
% ================================================================
\node[
    baseline,
    anchor=west,
    very thick,
    node font=\small\bfseries
] (mlp) at (10.05,2.68) {
    MLP baseline\\[2pt]
    $\widehat F=\widetilde F$
};

% ================================================================
% Soft-constrained branch
% ================================================================
\node[
    softouter,
    anchor=west,
    very thick
] (softbox) at (10.05,0.25) {};

\node[
    anchor=north,
    node font=\small\bfseries
] at ($(softbox.north)+(0,-0.17)$) {
    Soft MLP
};

\node[
    align=center
] (softraw) at ($(softbox.center)+(0,0.38)$) {
    $\widehat F=\widetilde F$
};

% ------------------------------------------------
% Movable soft-penalty box
% ------------------------------------------------
\def\softPenaltyXShift{0.0cm}   % + right, - left
\def\softPenaltyYShift{-0.5cm}  % + up,    - down

\node[
    innerlayer,
    fill=yellow!5,
    xshift=\softPenaltyXShift,
    yshift=\softPenaltyYShift
] (softpenalty) at (softbox.center) {
    \textbf{Physics penalty}\\[-0.1pt]
    added to training loss\\
    Eq.~(15)
};

% ================================================================
% Hard-constrained branch
% ================================================================
\node[
    hardouter,
    anchor=west,
    very thick,
    minimum height=96pt
] (hardbox) at (10.05,-3.26) {};

\node[
    anchor=north,
    node font=\small\bfseries
] at ($(hardbox.north)+(0,-0.16)$) {
    Ours (Hard-Constrained)
};

% Raw network output
\node[
    align=center
] (hardraw) at ($(hardbox.north)+(0,-0.78)$) {
    Raw output: $\widetilde F$
};

% ------------------------------------------------
% Movable enforcement layer
% ------------------------------------------------
\def\enforceXShift{0.0cm}   % + right, - left
\def\enforceYShift{-0.2cm}  % + up,    - down

\node[
    innerlayer,
    fill=green!6,
    xshift=\enforceXShift,
    yshift=\enforceYShift
] (enforce) at (hardbox.center) {
    \textbf{ReLU constraint-enforcement layer}\\[2pt]
    Pure lateral: Eq.~(10)\\
    Combined slip: Eq.~(11)
};

% Final feasible output
\node[
    align=center
] (hardout) at ($(hardbox.south)+(0,0.27)$) {
    $\widehat F\in\mathcal C^\tau(F_z,\mu)$
};

% ------------------------------------------------
% Internal hard-layer arrows
% ------------------------------------------------


% ================================================================
% Main signal flow
% ================================================================

% Common MLP -> baseline
\draw[
    conn,
    arrows=-Latex
]
    (shared.east)
    -- ++(0.60,0)
    |- (mlp.west);

% Common MLP -> soft branch
\draw[
    thick,
    arrows=-Latex
]
    (9.64,0)
    -- (10.07,0);

% Common MLP -> hard raw output
\draw[
    conn,
    arrows=-Latex
]
    (shared.east)
    -- ++(0.60,0)
    |- (10.07,-2.48);

% Task input -> shared MLP

% ================================================================
% Physical-context connections
% ================================================================

% Soft constraint: training-loss penalty
\draw[
    dashedconn
]
    (phys.east)
    -- ++(7,0)
    |- (softpenalty.west);

% Hard constraint: direct enforcement
\draw[
    conn,
    -
]
    (phys.east)
    -- ++(7,0)
    |- (10.07,-3.87);

\draw[
    ->,
    thick,
    dashed
]
    (10.07,-3.87)
    -- (10.63,-3.87);

  \draw[->, thick] (input.east) -- (shared.west);
\end{tikzpicture}
%
  }
  \caption{Hard-TireNet architecture.}
  \label{Tire:fig:architecture}
\end{figure}




 

\subsection{Force-Prediction Backbone}
\label{Tire:sec:backbone}

Let $\mv{h}_{F}^{(0)}:=\mv{x}$ be the combined-slip input.  A fully connected
force-prediction backbone with $L_F$ hidden layers is recursively defined by
\begin{equation}
  \label{Tire:eq:backbone:hidden}
  \mv{h}_{F}^{(\ell)}
  =
  \act_{\ell}
  \left(
    \wV_{F,\ell}^{\top}\mv{h}_{F}^{(\ell-1)}
  \right)
  ,
  \qquad \ell\in\{1,\ldots,L_F\}
  .
\end{equation}
The raw force output is
\begin{equation}
  \label{Tire:eq:backbone:output}
  \widehat{\mv{F}}_{u}
  =
  \wV_{F,L_F+1}^{\top}\mv{h}_{F}^{(L_F)}
  .
\end{equation}
Bias terms are included by augmenting each activation with a constant node,
as in the DNN definition in Chapter~\ref{CONAC:chap:hard:CONAC}.  The raw
output is not treated as the final force estimate when the physical constraint
is active.  The force-backbone weight vector is denoted by $\wth_F$.

\subsection{HardNet++ Differentiable Constraint Layer}
\label{Tire:sec:dpl}

The envelope parameters $\mv{\eta}$ are known a priori, so the constraint
function $c$ in \eqref{Tire:eq:friction:ellipse:rewritten} is used directly by
the projection layer.  The HardNet++ output is therefore
\begin{equation}
  \label{Tire:eq:hardnet:tire:output}
  \widehat{\mv{F}}
  =
  \Pi_{\mathrm{HN++}}
  \left(
    \widehat{\mv{F}}_{u};\mv{x},\mv{\eta}
  \right)
  .
\end{equation}
The map $\Pi_{\mathrm{HN++}}$ is evaluated in the forward pass.  Therefore,
the loss gradient is propagated through the constraint-enforcement
operation, and the backbone is trained using the constrained output.  This
is distinct from clipping a prediction after the network has been trained.

Let $\mv{y}^{(0)}:=\widehat{\mv{F}}_{u}$ and define the nonnegative violation
residual at iteration $k$ by
\begin{equation}
  \label{Tire:eq:dpl:residual}
  r^{(k)}
  :=
  \myReLU
  \left(
    c
    \left(
      \mv{y}^{(k)};\mv{x},\mv{\eta}
    \right)
  \right)
  .
\end{equation}
The constraint Jacobian is
\begin{equation}
  \label{Tire:eq:dpl:jacobian}
  \mm{J}^{(k)}
  :=
  \left.
  \ppfrac{
    c
    \left(
      \mv{y};\mv{x},\mv{\eta}
    \right)
  }{\mv{y}}
  \right|_{\mv{y}=\mv{y}^{(k)}}
  .
\end{equation}
Using a scalar damping parameter $\gamma_{\mathrm{dpl}}>0$, the differentiable
correction is
\begin{equation}
  \label{Tire:eq:dpl:update}
  \mv{y}^{(k+1)}
  =
  \mv{y}^{(k)}
  -
  {\mm{J}^{(k)}}^{\top}
  \left(
    \mm{J}^{(k)}{\mm{J}^{(k)}}^{\top}
    +\gamma_{\mathrm{dpl}}
  \right)^{-1}
  r^{(k)}
  .
\end{equation}
The iteration stops when $r^{(k)}\leq\varepsilon_{\mathrm{tol}}$ or when the
maximum number of iterations $K$ is reached.  If $k_*$ is the stopping index,
the final output is $\Pi_{\mathrm{HN++}}:=\mv{y}^{(k_*)}$.  The damping term
improves numerical conditioning when the constraint Jacobian is poorly
conditioned, which is particularly relevant near low-load or low-friction
operating points.

\begin{algorithm}[!t]
  \caption{HardNet++ combined-slip enforcement layer}
  \label{Tire:alg:dpl}
  \begin{algorithmic}[1]
    \STATE \textbf{Input:}
    Raw force $\widehat{\mv{F}}_u$, physical input $\mv{x}$,
    envelope parameters $\mv{\eta}$, damping parameter
    $\gamma_{\mathrm{dpl}}$, iteration limit $K$, and tolerance
    $\varepsilon_{\mathrm{tol}}$

    \STATE Initialize $\mv{y}^{(0)}\leftarrow\widehat{\mv{F}}_{u}$

    \FOR{$k=0$ to $K-1$}
      \STATE Evaluate
      $r^{(k)}\leftarrow
      \myReLU\left(
        c
        \left(
          \mv{y}^{(k)};\mv{x},\mv{\eta}
        \right)
      \right)$
      \IF{$|r^{(k)}|\leq\varepsilon_{\mathrm{tol}}$}
        \STATE \textbf{break}
      \ENDIF
      \STATE Evaluate the Jacobian $\mm{J}^{(k)}$
      \STATE Compute
      $\Delta\mv{y}^{(k)}\leftarrow
      -{\mm{J}^{(k)}}^{\top}
      (\mm{J}^{(k)}{\mm{J}^{(k)}}^{\top}+\gamma_{\mathrm{dpl}})^{-1}
      r^{(k)}$
      \STATE Update $\mv{y}^{(k+1)}\leftarrow
      \mv{y}^{(k)}+\Delta\mv{y}^{(k)}$
    \ENDFOR

    \STATE \textbf{Output:}
    $\widehat{\mv{F}}\leftarrow\mv{y}^{(k_*)}$
  \end{algorithmic}
\end{algorithm}

For the combined-slip ellipse, the constraint Jacobian is
\begin{equation}
  \label{Tire:eq:ellipse:jacobian}
  \ppfrac{c}{\mv{F}}
  =
  \begin{bmatrix}
    \dfrac{2F_x}{F_{x,\lim}(\mv{x};\mv{\eta})^2} &
    \dfrac{2F_y}{F_{y,\lim}(\mv{x};\mv{\eta})^2}
  \end{bmatrix}
  .
\end{equation}
When the DPL converges to a point for which the residual is below the
tolerance, the output satisfies the constraint up to that
tolerance:
\begin{equation}
  \label{Tire:eq:feasibility:tolerance}
  c
  \left(
      \widehat{\mv{F}};\mv{x},\mv{\eta}
  \right)
  \leq
  \varepsilon_{\mathrm{tol}}
  .
\end{equation}
This residual is a numerical enforcement error caused by a finite
iteration tolerance, since the envelope is known exactly.

%%%%%%%%%%%%%%%%%%%%%%%%%%%%%%%%%%%%%%%%%%%%%%%%%%%%%%%
% TRAINING AND ADAPTATION
%%%%%%%%%%%%%%%%%%%%%%%%%%%%%%%%%%%%%%%%%%%%%%%%%%%%%%%
\section{Training and Online Adaptation}
\label{Tire:sec:training}

\subsection{Training Objective}
\label{Tire:sec:training:objective}

Let a minibatch contain $B$ samples
$\{(\mv{x}_j,\mv{F}_j)\}_{j=1}^{B}$.  The constrained prediction for sample
$j$ is denoted by $\widehat{\mv{F}}_j$ and is obtained from
\eqref{Tire:eq:hardnet:tire:output}.  The force-fitting loss is
\begin{equation}
  \label{Tire:eq:fit:loss}
  \mathcal{L}_{\mathrm{fit}}
  :=
  \frac{1}{B}
  \sum_{j=1}^{B}
  \norm{
    \widehat{\mv{F}}_j-\mv{F}_j
  }^2
  .
\end{equation}
The constraint residual can be included as a diagnostic or as a numerical
regularizer:
\begin{equation}
  \label{Tire:eq:constraint:loss}
  \mathcal{L}_{\mathrm{con}}
  :=
  \frac{1}{B}
  \sum_{j=1}^{B}
  \left[
    \myReLU
    \left(
      c
      \left(
         \widehat{\mv{F}}_j;\mv{x}_j,\mv{\eta}
      \right)
    \right)
  \right]^2
  .
\end{equation}
The model is trained with
\begin{equation}
  \label{Tire:eq:total:loss}
  \mathcal{L}_{\mathrm{HN++}}
  =
  \mathcal{L}_{\mathrm{fit}}
  +
  \lambda_{\mathrm{con}}\mathcal{L}_{\mathrm{con}}
  +
  \lambda_{\mathrm{reg}}\mathcal{L}_{\mathrm{reg}}
  .
\end{equation}
Here, $\lambda_{\mathrm{con}},\lambda_{\mathrm{reg}}\geq0$ are
loss weights and $\mathcal{L}_{\mathrm{reg}}$ regularizes the network
parameters.  Because the forward output is already constrained,
$\mathcal{L}_{\mathrm{con}}$ is not the mechanism responsible for feasibility;
it is only a diagnostic or numerical regularizer.

\subsection{Online Parameter Adaptation}
\label{Tire:sec:online:adaptation}

For changing road and tire conditions, the weights can be adapted online.
Using the current minibatch loss, a projected gradient update is
\begin{equation}
  \label{Tire:eq:adaptation}
  \wth_F^{+}
  =
  \myproj_{\Theta_F}
  \left(
    \wth_F
    -\alpha_F
    \pptfrac{\mathcal{L}_{\mathrm{HN++}}}{\wth_F}
  \right)
  .
\end{equation}
Here, $\alpha_F\in\R_{>0}$ is a step size and the compact
set $\Theta_F$ bounds the parameters.  The
parameter projection prevents unbounded adaptation, while the HardNet++ layer
independently enforces the force constraint.

The online update is intentionally separated from output enforcement.  Even
if the weights are temporarily inaccurate, each forward pass
still evaluates the constraint layer with the known envelope.  Thus,
adaptation changes the backbone but does
not replace the hard projection step.

%%%%%%%%%%%%%%%%%%%%%%%%%%%%%%%%%%%%%%%%%%%%%%%%%%%%%%%
% VALIDATION
%%%%%%%%%%%%%%%%%%%%%%%%%%%%%%%%%%%%%%%%%%%%%%%%%%%%%%%
\section{Validation Protocol}
\label{Tire:sec:validation}

The validation study will focus exclusively on the combined-slip force map in
\eqref{Tire:eq:force:map}.  Simulation-generated tire-force data will be used
first so that the force labels and operating conditions are available.  The
benchmark will include multiple vertical-load and friction levels together
with combined variations of $\kappa$ and $\alpha$.  A subsequent study can
replace or augment the simulation data with measured tire data.

\subsection{Compared Models}
\label{Tire:sec:compared:models}

The following models will use the same backbone size, normalization,
train-validation split, optimizer, and stopping criterion:

\begin{enumerate}
  \item An unconstrained MLP that uses $\widehat{\mv{F}}_u$ directly as the
  predicted force.
  \item A soft-constrained MLP that uses the raw output and adds a finite
  physical penalty to the fitting loss.
  \item Hard-TireNet with the nonlinear differentiable enforcement layer in
  Algorithm~\ref{Tire:alg:dpl}.
\end{enumerate}

This comparison separates the inductive bias of a physical force envelope
from the difference between soft and hard enforcement.

\subsection{Operating Conditions}
\label{Tire:sec:operating:conditions}

The experiments will evaluate the following conditions:

\begin{enumerate}
  \item \textbf{IID:} random train-validation-test partitions over the
  available operating region.
  \item \textbf{Sparse:} a reduced training set with the validation and test
  sets unchanged.
  \item \textbf{Slip-OOD:} training and testing ranges are separated in
  $\kappa$ and/or $\alpha$.
  \item \textbf{Environment-OOD:} complete load-friction conditions are held
  out from training.
  \item \textbf{Sparse+OOD:} sparse training is combined with slip and
  environment shifts.
\end{enumerate}

All normalization statistics and envelope-calibration parameters will be
computed using the training partition only.  This avoids transferring
information from the test distribution into the physical constraint model.

\subsection{Evaluation Metrics}
\label{Tire:sec:metrics}

Prediction accuracy will be quantified using component-wise root
mean-square error:
\begin{equation}
  \label{Tire:eq:rmse}
  \mathrm{RMSE}_{i}
  :=
  \sqrt{
    \frac{1}{N_{\mathrm{test}}}
    \sum_{j=1}^{N_{\mathrm{test}}}
    \left(
      \widehat{F}_{i,j}-F_{i,j}
    \right)^2
  }
  ,
  \qquad
  i\in\{x,y\}
  .
\end{equation}
The element-wise force error is also reported as
\begin{equation}
  \label{Tire:eq:rmse:element}
  \mathrm{RMSE}_{\mathrm{elem}}
  :=
  \sqrt{
    \frac{1}{2N_{\mathrm{test}}}
    \sum_{j=1}^{N_{\mathrm{test}}}
    \norm{
      \widehat{\mv{F}}_j-\mv{F}_j
    }^2
  }
  .
\end{equation}
Physical consistency will be quantified by the constraint-violation rate,
computed using the same envelope used during inference:
\begin{equation}
  \label{Tire:eq:violation:rate}
  r_{\mathrm{viol}}
  :=
  \frac{100}{N_{\mathrm{test}}}
  \sum_{j=1}^{N_{\mathrm{test}}}
  \mathbb{I}
  \left[
    c
    \left(
      \widehat{\mv{F}}_j;\mv{x}_j,\mv{\eta}_j
    \right)
    >
    \varepsilon_{\mathrm{tol}}
  \right]
  .
\end{equation}
The enforcement residual, the number of DPL iterations, and the inference
time will also be recorded.  The central feasibility target is
$r_{\mathrm{viol}}=0$ with respect to the known envelope, while
accuracy and computational cost determine whether the constrained surrogate
is suitable for real-time vehicle estimation or model predictive control.

%%%%%%%%%%%%%%%%%%%%%%%%%%%%%%%%%%%%%%%%%%%%%%%%%%%%%%%
% CONCLUSION
%%%%%%%%%%%%%%%%%%%%%%%%%%%%%%%%%%%%%%%%%%%%%%%%%%%%%%%
\section{Conclusion}
\label{Tire:sec:conclusion}

This chapter proposed Hard-TireNet, a hard-constrained neural surrogate
for combined longitudinal and lateral tire-force estimation.  The method
uses a neural backbone to approximate the tire-force map and a
HardNet++ differentiable projection layer to enforce the known
nonlinear friction-ellipse constraint.

The proposed formulation provides a direct bridge between the physical
tire-force problem and the general hard-constraint framework developed in
this proposal.  Its main hypothesis is that hard enforcement will be most
useful under sparse data and distribution shift, where an unconstrained
network can extrapolate to force combinations that are inconsistent with the
available friction.  The planned experiments will test this hypothesis
through prediction accuracy, violation rate, numerical residual, and
real-time computational cost.

The initial version assumes that $F_z$ and $\mu$ are available and that a
conservative envelope can be calibrated.  Future work will address
uncertainty in these quantities, online identification of the admissible set,
measured tire data, and integration of the constrained surrogate into a
vehicle-level observer or predictive controller.
